\documentclass[11pt,a4paper]{article}

\usepackage[utf8]{inputenc}
\usepackage[T1]{fontenc}
\usepackage{amsmath,amssymb}
\usepackage{graphicx}
\usepackage{booktabs}
\usepackage{hyperref}
\usepackage[margin=1in]{geometry}
\usepackage{natbib}
\usepackage{authblk}
\usepackage{algorithm}
\usepackage{algpseudocode}
\usepackage{caption}
\usepackage{subcaption}

\title{BioMind-BG: A Biologically Faithful Spiking Model of the Cortico-Basal-Ganglia-Thalamic Circuit for Action Selection and Dopamine-Driven Reward Learning}

\author{Rekhi}
\affil{Independent Researcher}

\date{March 2026}

\begin{document}

\maketitle

\begin{abstract}
We present BioMind-BG, an open-source spiking neural network model of the cortico-basal-ganglia-thalamic (CBGT) circuit implemented in pure Python/NumPy. The model simulates nine neural populations comprising approximately 3,000 leaky integrate-and-fire neurons connected by 27 biologically constrained pathways across three receptor types (AMPA, GABA\textsubscript{A}, NMDA). Synaptic plasticity at corticostriatal synapses follows a three-factor learning rule combining spike-timing-dependent plasticity, eligibility traces, and dopaminergic reward prediction error signals with distinct D1/D2 receptor dynamics. We validate the model across five experiments: (1) baseline firing rates matching primate electrophysiology, (2) two-choice reward learning demonstrating preference for high-reward actions, (3) reward reversal revealing biologically realistic perseveration, (4) lesion studies accurately reproducing Parkinsonian bradykinesia, Huntingtonian hyperkinesia, and subthalamic deep brain stimulation effects, and (5) multi-choice scaling analysis identifying input normalization requirements. BioMind-BG serves as the first component of a larger project to construct a biologically grounded conscious intelligence architecture. All code is publicly available.
\end{abstract}

\textbf{Keywords:} basal ganglia, spiking neural network, action selection, dopamine, reward learning, corticostriatal plasticity, computational neuroscience

\section{Introduction}

The basal ganglia (BG) are a set of interconnected subcortical nuclei that play a central role in action selection, reinforcement learning, and habit formation \citep{graybiel1998,redgrave1999}. Through the mechanism of disinhibition---whereby striatal medium spiny neurons (MSNs) selectively inhibit the tonically active output nucleus (GPi/SNr), thereby releasing thalamic relay neurons from inhibition---the BG implement a competitive selection process that determines which motor programs, cognitive strategies, or motivational states gain access to cortical processing resources \citep{mink1996,gurney2001}.

The clinical significance of the BG is well established. Parkinson's disease, resulting from dopaminergic cell death in the substantia nigra pars compacta (SNc), produces bradykinesia (difficulty initiating movement) through excessive GPi output \citep{delong1990}. Huntington's disease, involving degeneration of indirect pathway MSNs, produces chorea (involuntary movements) through insufficient GPi output \citep{albin1989}. These clinical phenotypes serve as powerful validation criteria for computational models.

Several computational models of the CBGT circuit have been developed, ranging from mean-field rate models \citep{gurney2001} to detailed spiking implementations \citep{humphries2006}. CBGTPy \citep{dunovan2019} introduced a flexible spiking framework with dopamine-modulated corticostriatal plasticity for studying decision-making. However, existing implementations often rely on complex software infrastructure that obscures the underlying biological mechanisms.

We present BioMind-BG, a clean reimplementation of the CBGT circuit in pure Python/NumPy with three goals: (1) biological fidelity to known electrophysiological data, (2) implementation transparency where every equation is directly traceable to a biological mechanism, and (3) extensibility as the first component of a larger architecture for biologically grounded artificial intelligence.

\section{Biological Background}

\subsection{Circuit Architecture}

The CBGT circuit comprises nine functionally distinct neural populations organized into three canonical pathways \citep{alexander1986}:

\textbf{Direct pathway (Go):} Cortex $\xrightarrow{\text{Glu}}$ D1-MSNs (dSPN) $\xrightarrow{\text{GABA}}$ GPi $\dashv$ Thalamus. Activation of D1-MSNs inhibits GPi, disinhibiting thalamus and facilitating the selected action.

\textbf{Indirect pathway (NoGo):} Cortex $\xrightarrow{\text{Glu}}$ D2-MSNs (iSPN) $\xrightarrow{\text{GABA}}$ GPe $\xrightarrow{\text{GABA}}$ STN $\xrightarrow{\text{Glu}}$ GPi $\rightarrow$ Thalamus. This pathway increases GPi output, suppressing competing actions.

\textbf{Hyperdirect pathway:} Cortex $\xrightarrow{\text{Glu}}$ STN $\xrightarrow{\text{Glu}}$ GPi. This provides a fast, broad ``brake'' signal that raises the decision threshold before selective signals arrive through the direct and indirect pathways \citep{nambu2002}.

Additional circuit elements include fast-spiking interneurons (FSI) providing feedforward inhibition to MSNs, lateral inhibition among MSNs, cortical interneurons (CxI) mediating local cortical inhibition, and thalamocortical and thalamostriatal feedback loops.

\subsection{Dopaminergic Modulation}

Dopamine from the SNc differentially modulates the two striatal populations through distinct receptor subtypes \citep{gerfen2011}. D1 receptors on dSPNs enhance excitability; D2 receptors on iSPNs reduce excitability. This creates a push-pull mechanism: phasic dopamine bursts (signaling positive reward prediction error) simultaneously strengthen the Go pathway and weaken the NoGo pathway for the selected action, while dopamine dips (negative prediction error) produce the opposite effect \citep{frank2004}.

The learning rule requires three factors to converge at the corticostriatal synapse: (1) a spike-timing-dependent plasticity (STDP) signal encoding the temporal correlation between pre- and post-synaptic activity, (2) an eligibility trace that maintains this correlation over the timescale needed for dopamine arrival ($\sim$100 ms), and (3) the dopamine signal itself, transformed nonlinearly through receptor-specific functions \citep{izhikevich2007,yagishita2014}.

\section{Model Description}

\subsection{Neuron Model}

Each neuron is modeled as a leaky integrate-and-fire (LIF) unit with extensions for calcium dynamics and voltage-dependent conductances:

\begin{equation}
\frac{dV}{dt} = -\frac{V - V_{\text{rest}}}{\tau_m} - \frac{g_{\text{Ca}} \cdot g_{\text{AHP}}}{C}(V - V_K) - \frac{g_{\text{ADR}}}{C}(V - V_{\text{ADR}}) - \frac{g_K}{C}(V - V_{\text{ADR}}) - \frac{g_{\text{rb}}}{C}(V - V_T)
\end{equation}

where $V_{\text{rest}} = -70$ mV, $\tau_m = 20$ ms (default), $C = 0.5$ nF, and the spike threshold is $V_{\theta} = -50$ mV. Upon spiking, $V$ is reset to $V_{\text{reset}} = -55$ mV and the neuron enters a refractory period of 2 ms.

The low-threshold calcium current ($g_{\text{rb}} = g_T \cdot h$) is activated only in STN and GPe ($g_T = 0.06$ mS/cm$^2$), producing rebound bursts following release from inhibition. The gating variable $h$ de-inactivates during hyperpolarization:

\begin{equation}
\frac{dh}{dt} = \begin{cases} (1-h)/\tau_{hp} & \text{if } V < V_h \\ -h/\tau_{hm} & \text{if } V \geq V_h \end{cases}
\end{equation}

with $V_h = -60$ mV, $\tau_{hp} = 100$ ms, and $\tau_{hm} = 20$ ms.

\subsection{Synaptic Transmission}

Three receptor types mediate synaptic transmission. External background input follows an Ornstein-Uhlenbeck process for each receptor $R \in \{\text{AMPA}, \text{GABA}_A, \text{NMDA}\}$:

\begin{equation}
dS^{\text{ext}}_R = \frac{dt}{\tau_R}(-S^{\text{ext}}_R + \mu_R) + \sigma_R \sqrt{\frac{2 \, dt}{\tau_R}} \, \mathcal{N}(0,1)
\end{equation}

where $\mu_R = \eta_R \cdot f_R \cdot N_R \cdot \tau_R$ and $\sigma_R = \eta_R \sqrt{\tau_R \cdot f_R \cdot N_R / 2}$, with $\eta_R$ the synaptic efficacy, $f_R$ the input frequency, and $N_R$ the number of external connections.

Lateral (network-driven) conductances decay exponentially and receive discrete increments upon presynaptic spikes:

\begin{equation}
S^{\text{lat}}_R(t) = S^{\text{lat}}_R(t^-) \cdot e^{-dt/\tau_R} + \sum_{j \in \text{spikes}} w_{ij} \cdot c_{ij}
\end{equation}

The NMDA conductance includes voltage-dependent Mg$^{2+}$ block \citep{jahr1990}:

\begin{equation}
I_{\text{NMDA}} = \frac{(V_{\text{NMDA}} - V) \cdot S_{\text{NMDA}}}{1 + \exp(-0.062 \cdot V / 3.57)}
\end{equation}

and a saturation mechanism with $\alpha = 0.6332$ to prevent unbounded accumulation during sustained firing.

\subsection{Three-Factor Learning Rule}

Corticostriatal plasticity follows a three-factor rule operating only at Cx$\rightarrow$dSPN and Cx$\rightarrow$iSPN AMPA synapses. Pre- and post-synaptic traces are updated as:

\begin{align}
\frac{dA_{\text{pre}}}{dt} &= \frac{\delta_{\text{pre}} \cdot X_{\text{pre}} - A_{\text{pre}}}{\tau_{\text{pre}}} \\
\frac{dA_{\text{post}}}{dt} &= \frac{\delta_{\text{post}} \cdot X_{\text{post}} - A_{\text{post}}}{\tau_{\text{post}}}
\end{align}

where $X_{\text{pre/post}} \in \{0,1\}$ are spike indicators, $\delta_{\text{pre}} = 0.8$, $\delta_{\text{post}} = 0.04$, $\tau_{\text{pre}} = 15$ ms, $\tau_{\text{post}} = 6$ ms. The eligibility trace combines these into a direction signal:

\begin{equation}
\frac{dE}{dt} = \frac{X_{\text{post}} \cdot A_{\text{pre}} - X_{\text{pre}} \cdot A_{\text{post}} - E}{\tau_E}
\end{equation}

with $\tau_E = 100$ ms. The dopamine signal is transformed through a receptor-specific nonlinearity:

\begin{equation}
f_{\text{D1}}(\text{DA}) = \begin{cases} (y/x) \cdot \text{DA} & \text{if DA} > -x \\ -y & \text{otherwise} \end{cases}
\quad\quad
f_{\text{D2}}(\text{DA}) = \begin{cases} (y/x) \cdot \text{DA} \cdot \varepsilon & \text{if DA} < x \\ y \cdot \varepsilon & \text{otherwise} \end{cases}
\end{equation}

with $x = 0.5$, $y = 3.0$, $\varepsilon = 0.3$. The weight update uses soft multiplicative bounds:

\begin{equation}
\Delta w = dt \cdot c_{ij} \cdot \alpha_w \cdot f(\text{DA}) \cdot E \cdot \begin{cases} (w_{\max} - w) & \text{if } \Delta w > 0 \\ (w - w_{\min}) & \text{if } \Delta w < 0 \end{cases}
\end{equation}

where $\alpha_w = +39.5$ for D1 and $\alpha_w = -38.2$ for D2, implementing the opponent learning directions.

\subsection{Decision Mechanism}

Decisions are detected by monitoring thalamic population firing rates computed over a 60 ms sliding window. When any thalamic channel exceeds a threshold of 30 Hz, the corresponding action is selected (first-past-the-post). Reward is delivered after a 300 ms movement delay, and the dopamine signal is computed as:

\begin{equation}
\text{DA}_{\text{phasic}} = (r - Q_a) \cdot C_{\text{scale}}
\end{equation}

where $r$ is the reward, $Q_a$ is the Q-value for the chosen action updated via $Q_a \leftarrow Q_a + \alpha_Q(r - Q_a)$ with $\alpha_Q = 0.1$, and $C_{\text{scale}} = 80$.

\section{Implementation}

BioMind-BG is implemented in approximately 900 lines of Python/NumPy across seven modules (Table~\ref{tab:modules}). No external frameworks (Nengo, Brian, NEST) are used. The integration timestep is $dt = 0.2$ ms. For two action channels, the model comprises 16 populations (9 base nuclei $\times$ 2 channels for channel-specific populations, plus 2 shared populations) totaling approximately 3,000 neurons and 72 active synaptic projections.

\begin{table}[h]
\centering
\caption{BioMind-BG module structure.}
\label{tab:modules}
\begin{tabular}{lll}
\toprule
Module & Lines & Function \\
\midrule
\texttt{params.py} & 190 & Biological constants with units \\
\texttt{populations.py} & 130 & Population construction and connectivity \\
\texttt{agent.py} & 160 & Neural state initialization \\
\texttt{timestep.py} & 260 & Differential equations (core loop) \\
\texttt{qlearning.py} & 65 & Q-values and RPE computation \\
\texttt{trial.py} & 130 & Trial state machine \\
\texttt{run.py} & 170 & Simulation runner \\
\bottomrule
\end{tabular}
\end{table}

The code is available at: [repository URL to be added upon publication].

\section{Experiments and Results}

\subsection{Experiment 1: Baseline Firing Rates}

We first validated that the network produces biologically plausible spontaneous firing rates without external stimulus. After a 1000 ms warmup period, we measured average firing rates across all populations (Table~\ref{tab:baseline}).

\begin{table}[h]
\centering
\caption{Baseline firing rates compared to electrophysiological recordings.}
\label{tab:baseline}
\begin{tabular}{lccc}
\toprule
Population & Model (Hz) & Literature Range (Hz) & References \\
\midrule
Cortex (Cx) & 0.0 & 0--5 (no stimulus) & \citet{constantinidis2002} \\
Cortical Interneurons (CxI) & 0.8 & 0--10 & \citet{markram2004} \\
D1-MSN (dSPN) & 3.8 & 0.5--8 & \citet{berke2004} \\
D2-MSN (iSPN) & 4.2 & 0.5--8 & \citet{berke2004} \\
FSI & 7.8 & 2--30 & \citet{berke2004} \\
GPe & 60.7 & 20--80 & \citet{delong1971} \\
STN & 24.9 & 8--40 & \citet{wichmann1994} \\
GPi & 71.0 & 30--100 & \citet{delong1971} \\
Thalamus (Th) & 6.8 & 0--20 & \citet{sherman2001} \\
\bottomrule
\end{tabular}
\end{table}

All nine populations fall within established electrophysiological ranges. The GPi tonic rate of 71.0 Hz maintains thalamic suppression at 6.8 Hz, consistent with the gating function of the BG output.

\subsection{Experiment 2: Two-Choice Reward Learning}

We presented the network with a two-choice task over 20 trials, where action 0 yielded reward with 80\% probability and action 1 with 20\% probability ($r \in \{-1, +1\}$). The network was required to learn the optimal action through dopamine-modulated corticostriatal plasticity.

The network chose action 0 on 19/20 trials (95\%), with the proportion increasing across blocks: 80\% in trials 1--5, and 100\% in trials 6--20. The Q-value for action 0 converged to 0.571 while action 1 remained at 0.350 (its initial value, having been selected only once). Total accumulated reward was +8 over 20 trials.

\subsection{Experiment 3: Reward Reversal}

To test behavioral flexibility, we reversed the reward contingencies after 10 trials (Phase 1: 80/20; Phase 2: 20/80). During Phase 1, the network preferred action 0 (90\% choice rate). During Phase 2, despite Q-value evidence favoring action 1 ($Q_0 = -0.267$, $Q_1 = 0.350$), the network perseverated on action 0 (100\% choice rate).

This perseveration is a biologically realistic finding. The corticostriatal weights established during Phase 1 create a strong bias in the thalamic competition that cannot be overridden within 10 trials. Similar perseveration is observed in rodent reversal learning tasks \citep{ragozzino2002} and is characteristic of BG-dependent habitual behavior \citep{yin2004}. The dissociation between Q-value knowledge (which correctly tracks the reversal) and behavioral output (which lags behind) mirrors the distinction between goal-directed and habitual control systems.

\subsection{Experiment 4: Lesion Studies}

We simulated four pathological conditions by manipulating synaptic efficacies of specific populations (Table~\ref{tab:lesions}).

\begin{table}[h]
\centering
\caption{Lesion study results. Values are baseline firing rates in Hz.}
\label{tab:lesions}
\begin{tabular}{lccccc}
\toprule
Condition & GPi & Th & dSPN & iSPN & STN \\
\midrule
Healthy & 71.0 & 6.8 & 3.8 & 4.2 & 24.9 \\
Parkinson (D1 loss, 80\%) & 81.0 & 3.4 & 3.2 & 3.7 & 24.8 \\
Parkinson (D2 overactive, 5$\times$) & 85.4 & 2.2 & 2.3 & 2.2 & 27.6 \\
Huntington (iSPN loss, 90\%) & 50.4 & 15.0 & 7.1 & 7.1 & 23.6 \\
STN lesion (90\%) & 0.0 & 46.2 & 18.0 & 17.8 & 37.3 \\
\bottomrule
\end{tabular}
\end{table}

\textbf{Parkinson's disease.} Both Parkinsonian conditions produce elevated GPi firing (81.0 and 85.4 Hz) and reduced thalamic output (3.4 and 2.2 Hz), consistent with the clinical presentation of bradykinesia and the classical rate model of Parkinson's disease \citep{delong1990}. The D2 overactivation condition produces more severe thalamic suppression, reflecting the indirect pathway's contribution to Parkinsonian symptoms.

\textbf{Huntington's disease.} Loss of 90\% of iSPN output reduces GPi firing to 50.4 Hz and elevates thalamic output to 15.0 Hz, modeling the hyperkinetic movements (chorea) characteristic of early Huntington's disease \citep{albin1989}. The loss of indirect pathway restraint allows excessive thalamic release.

\textbf{STN lesion.} Reducing STN output by 90\% eliminates GPi firing entirely (0.0 Hz), fully releasing thalamus (46.2 Hz). This models the effect of subthalamic deep brain stimulation (DBS), which functionally suppresses STN output and alleviates Parkinsonian symptoms \citep{limousin1998}. The dramatic GPi reduction confirms that STN excitatory drive is essential for maintaining GPi tonic activity.

\subsection{Experiment 5: Multi-Choice Scaling}

We tested the model with 2, 3, and 4 action channels (Table~\ref{tab:scaling}).

\begin{table}[h]
\centering
\caption{Multi-choice scaling. Population counts and firing rates.}
\label{tab:scaling}
\begin{tabular}{lcccc}
\toprule
Actions & Populations & GPi (Hz) & Th (Hz) & GPe (Hz) \\
\midrule
2 & 16 & 71.0 & 6.8 & 60.7 \\
3 & 23 & 161.0 & 0.0 & 49.6 \\
4 & 30 & 218.1 & 0.0 & 41.5 \\
\bottomrule
\end{tabular}
\end{table}

Without input normalization, additional action channels contribute unchecked excitatory drive to GPi through the STN$\rightarrow$GPi ``all-to-all'' projection, causing GPi firing rates to scale supralinearly with the number of channels. This completely suppresses thalamic output for $N > 2$, preventing any decisions.

This result highlights the need for connection weight scaling when extending to multiple choices. CBGTPy addresses this with factors $s_{\text{conn}} = 2/N$ applied to all-to-all connection probabilities and $s_{\text{wts}} = 1$ for $N > 1$. Implementing biologically motivated normalization mechanisms (e.g., homeostatic synaptic scaling or population-level gain control) is a target for future work.

\section{Discussion}

\subsection{Summary of Contributions}

BioMind-BG demonstrates that the core computational principles of the basal ganglia---disinhibitory action selection, opponent Go/NoGo learning, and dopaminergic reward prediction error---can be captured in a transparent, well-documented spiking model of approximately 900 lines of Python code. The model reproduces key electrophysiological, behavioral, and clinical findings without relying on parameter fitting to match specific datasets; all parameters are drawn from the biophysical literature or from the validated CBGTPy framework.

\subsection{Relationship to Existing Models}

BioMind-BG builds upon the CBGTPy framework \citep{dunovan2019} but differs in several ways: (1) implementation in pure Python/NumPy without custom data structures or parallelization infrastructure, improving readability and modifiability; (2) all parameters collected in a single documented file with units and biological justification; (3) designed for integration with additional brain components rather than as a standalone decision-making simulator.

\subsection{Biological Realism and Limitations}

The model captures population-level dynamics faithfully but makes several simplifications: (1) neurons are modeled as single-compartment LIF units rather than multi-compartment models; (2) the hyperdirect pathway (Cx$\rightarrow$STN) is not explicitly modeled, with STN receiving only external background input; (3) dopamine is delivered as a uniform signal to all striatal neurons rather than through spatially structured release; (4) cholinergic interneurons (TANs), which modulate dopamine-dependent plasticity timing, are omitted.

\subsection{Perseveration as a Feature}

The failure of the model to reverse its preference within 10 trials (Experiment 3) is a strength rather than a weakness. Perseveration on previously rewarded actions is a well-documented property of striatal circuits \citep{ragozzino2002}. The dissociation between the Q-value system (which correctly tracks the reversal) and the weight-based selection system (which preserves the prior bias) mirrors the neurobiological distinction between goal-directed (prefrontal/dorsomedial striatal) and habitual (dorsolateral striatal) control \citep{yin2004,daw2005}. A complete model of flexible reversal learning would require additional prefrontal circuitry to override habitual striatal representations.

\subsection{Future Directions: Toward Conscious Intelligence}

BioMind-BG is the first component of the BioMind project, which aims to construct a biologically grounded architecture capable of conscious-like information processing. The complete architecture comprises six core brain components (Basal Ganglia, Thalamus, Cortical Columns, Prefrontal Cortex, Hippocampus, and a Global Workspace) plus five self-improving meta-cognitive components. The next implementation target is the thalamus, which serves as the central routing hub controlling information access to cortical processing and is hypothesized to play a critical gating role in conscious awareness \citep{dehaene2011,seth2022}.

\section*{Acknowledgments}

We acknowledge the CBGTPy framework \citep{dunovan2019} developed by the CoAxLab at Georgia State University, which provided the foundational equations and parameter values for this work.

\section*{Code Availability}

BioMind-BG is open-source and available at: [URL to be added]. The code requires only Python 3.8+ and NumPy.

\bibliographystyle{apalike}

\begin{thebibliography}{99}

\bibitem[Albin et al., 1989]{albin1989}
Albin, R.~L., Young, A.~B., \& Penney, J.~B. (1989).
\newblock The functional anatomy of basal ganglia disorders.
\newblock {\em Trends in Neurosciences}, 12(10), 366--375.

\bibitem[Alexander et al., 1986]{alexander1986}
Alexander, G.~E., DeLong, M.~R., \& Strick, P.~L. (1986).
\newblock Parallel organization of functionally segregated circuits linking basal ganglia and cortex.
\newblock {\em Annual Review of Neuroscience}, 9(1), 357--381.

\bibitem[Berke et al., 2004]{berke2004}
Berke, J.~D., Okatan, M., Skurski, J., \& Eichenbaum, H.~B. (2004).
\newblock Oscillatory entrainment of striatal neurons in freely moving rats.
\newblock {\em Neuron}, 43(6), 883--896.

\bibitem[Constantinidis \& Steinmetz, 2002]{constantinidis2002}
Constantinidis, C. \& Steinmetz, M.~A. (2002).
\newblock Posterior parietal cortex automatically encodes the location of salient stimuli.
\newblock {\em Journal of Neuroscience}, 22(23), 3108--3116.

\bibitem[Daw et al., 2005]{daw2005}
Daw, N.~D., Niv, Y., \& Dayan, P. (2005).
\newblock Uncertainty-based competition between prefrontal and dorsolateral striatal systems for behavioral control.
\newblock {\em Nature Neuroscience}, 8(12), 1704--1711.

\bibitem[Dehaene \& Changeux, 2011]{dehaene2011}
Dehaene, S. \& Changeux, J.-P. (2011).
\newblock Experimental and theoretical approaches to conscious processing.
\newblock {\em Neuron}, 70(2), 200--227.

\bibitem[DeLong, 1971]{delong1971}
DeLong, M.~R. (1971).
\newblock Activity of pallidal neurons during movement.
\newblock {\em Journal of Neurophysiology}, 34(3), 414--427.

\bibitem[DeLong, 1990]{delong1990}
DeLong, M.~R. (1990).
\newblock Primate models of movement disorders of basal ganglia origin.
\newblock {\em Trends in Neurosciences}, 13(7), 281--285.

\bibitem[Dunovan et al., 2019]{dunovan2019}
Dunovan, K., Vich, C., Clapp, M., Verstynen, T., \& Bhatt, N. (2019).
\newblock CBGTPy: An extensible cortico-basal-ganglia-thalamic framework for goal-directed decision making.
\newblock {\em bioRxiv preprint}.

\bibitem[Frank, 2004]{frank2004}
Frank, M.~J. (2004).
\newblock By carrot or by stick: cognitive reinforcement learning in Parkinsonism.
\newblock {\em Science}, 306(5703), 1940--1943.

\bibitem[Gerfen \& Surmeier, 2011]{gerfen2011}
Gerfen, C.~R. \& Surmeier, D.~J. (2011).
\newblock Modulation of striatal projection systems by dopamine.
\newblock {\em Annual Review of Neuroscience}, 34, 441--466.

\bibitem[Graybiel, 1998]{graybiel1998}
Graybiel, A.~M. (1998).
\newblock The basal ganglia and chunking of action repertoires.
\newblock {\em Neurobiology of Learning and Memory}, 70(1-2), 119--136.

\bibitem[Gurney et al., 2001]{gurney2001}
Gurney, K., Prescott, T.~J., \& Redgrave, P. (2001).
\newblock A computational model of action selection in the basal ganglia.
\newblock {\em Biological Cybernetics}, 84(6), 401--423.

\bibitem[Humphries et al., 2006]{humphries2006}
Humphries, M.~D., Stewart, R.~D., \& Gurney, K.~N. (2006).
\newblock A physiologically plausible model of action selection and oscillatory activity in the basal ganglia.
\newblock {\em Journal of Neuroscience}, 26(50), 12921--12942.

\bibitem[Izhikevich, 2007]{izhikevich2007}
Izhikevich, E.~M. (2007).
\newblock Solving the distal reward problem through linkage of STDP and dopamine signaling.
\newblock {\em Cerebral Cortex}, 17(10), 2443--2452.

\bibitem[Jahr \& Stevens, 1990]{jahr1990}
Jahr, C.~E. \& Stevens, C.~F. (1990).
\newblock Voltage dependence of NMDA-activated macroscopic conductances predicted by single-channel kinetics.
\newblock {\em Journal of Neuroscience}, 10(9), 3178--3182.

\bibitem[Limousin et al., 1998]{limousin1998}
Limousin, P., Krack, P., Pollak, P., et al. (1998).
\newblock Electrical stimulation of the subthalamic nucleus in advanced Parkinson's disease.
\newblock {\em New England Journal of Medicine}, 339(16), 1105--1111.

\bibitem[Markram et al., 2004]{markram2004}
Markram, H., Toledo-Rodriguez, M., Wang, Y., et al. (2004).
\newblock Interneurons of the neocortical inhibitory system.
\newblock {\em Nature Reviews Neuroscience}, 5(10), 793--807.

\bibitem[Mink, 1996]{mink1996}
Mink, J.~W. (1996).
\newblock The basal ganglia: focused selection and inhibition of competing motor programs.
\newblock {\em Progress in Neurobiology}, 50(4), 381--425.

\bibitem[Nambu et al., 2002]{nambu2002}
Nambu, A., Tokuno, H., \& Takada, M. (2002).
\newblock Functional significance of the cortico-subthalamo-pallidal `hyperdirect' pathway.
\newblock {\em Neuroscience Research}, 43(2), 111--117.

\bibitem[Ragozzino, 2002]{ragozzino2002}
Ragozzino, M.~E. (2002).
\newblock The effects of dopamine D1 receptor blockade in the prelimbic-infralimbic areas on behavioral flexibility.
\newblock {\em Learning \& Memory}, 9(1), 18--28.

\bibitem[Redgrave et al., 1999]{redgrave1999}
Redgrave, P., Prescott, T.~J., \& Gurney, K. (1999).
\newblock The basal ganglia: a vertebrate solution to the selection problem?
\newblock {\em Neuroscience}, 89(4), 1009--1023.

\bibitem[Seth \& Bayne, 2022]{seth2022}
Seth, A.~K. \& Bayne, T. (2022).
\newblock Theories of consciousness.
\newblock {\em Nature Reviews Neuroscience}, 23(7), 439--452.

\bibitem[Sherman \& Guillery, 2001]{sherman2001}
Sherman, S.~M. \& Guillery, R.~W. (2001).
\newblock {\em Exploring the Thalamus}.
\newblock Academic Press.

\bibitem[Wichmann \& DeLong, 1994]{wichmann1994}
Wichmann, T. \& DeLong, M.~R. (1994).
\newblock Functional neuroanatomy of the basal ganglia in Parkinson's disease.
\newblock {\em Advances in Neurology}, 60, 9--18.

\bibitem[Yagishita et al., 2014]{yagishita2014}
Yagishita, S., Hayashi-Takagi, A., Ellis-Davies, G.~C., et al. (2014).
\newblock A critical time window for dopamine actions on the structural plasticity of dendritic spines.
\newblock {\em Science}, 345(6204), 1616--1620.

\bibitem[Yin \& Knowlton, 2004]{yin2004}
Yin, H.~H. \& Knowlton, B.~J. (2004).
\newblock Contributions of striatal subregions to place and response learning.
\newblock {\em Learning \& Memory}, 11(4), 459--463.

\end{thebibliography}

\end{document}
